% !TEX root = ../main.tex
\lecture{6}{2026-10-26}{Unternehmensbewertung}

\section{Unternehmensbewertung}
Nicht den richtigen Wert eines Unternehmens berechnen, sondern Plausibilität des Aktienpreises über Annahmen überprüfen.
Reverse Valuation
\subsection{Einführung}

\subsubsection{Grundsatzfragen}
\begin{itemize}
    \item Investitionsrechnung
    \item Unternehmensbewertung
    \item Grundprobleme
    \begin{itemize}
        \item ''Unternehmenswert ist \textit{subjektiv}'' 
        \item Unternehmens ist meist grösser als die Summe der Vermögenswerte \footnote{Ja klar, es wird nicht nach dem Wert der einzelnen Waren/Geld/Anleihen gesucht, die das Kapital ausmachen, sondern den Wert des Kapitals selbst gesucht. Das heisst, welche Geldgrösse zu einer Profitrate ist im Stande eine G}
    \end{itemize}
\end{itemize}
\subsection{Substanzwert-Methode}
\begin{itemize}
    \item Entspricht der Summe aller Vermögenswerte
    \begin{itemize}
        \item Aus Infos der Bilanzrechnungen berechnet
        \item Gegenwartbezogen, entspricht nicht ganz den zukünftigen Wert
        \item objektive Bilanzbetrachtung
    \end{itemize}
    \item Berechnung:
    \begin{itemize}
        \item Netto Substanzwert
        \begin{itemize}
            \item Wie viel man für das Eigenkapital anteil zahlen muss
            \item UV + AV - FK = EK oder Netto Substanzwert
        \end{itemize}
    \end{itemize}
\end{itemize}

\subsubsection{Stillen Reserven}
\begin{definition}
    [Stillen Reserven] Differenz zwischen Bilanzwert der Vermögensteile (Buchwert aus Finanzbuchhaltung) und tatsächlichem Wert (Börse/Marktwert) \footnote{Werden hier nicht qualitativ verschiedene Sachen in Geld verglichen?}
\end{definition}
\begin{itemize}
    \item Müssen Sauber berechnet werden, um Gläubiger zu schützen und Transparenz bewahren
\end{itemize}

\subsubsection{Bereinigung}
\begin{itemize}
    \item 
\end{itemize}

\subsection{Ertragswert}

\begin{definition}
    [Kapitalkosten] sind Werte/Rendite, die an Kapitalgeber bezahlt werden müssen relativ zum gegebenem Gesamtkapital.
\end{definition}
\textbf{Kapitalkostenarten}
\begin{itemize}
    \item Fremdkapitalkosten
    \begin{itemize}
        \item Explizit vereinbarte \textit{Zins}zahlungen
    \end{itemize}
    \item Eigenkapitalkosten
    \begin{itemize}
        \item Implizite Kapitalkosten als Opportunitätskosten. Eigentümer erwarten Redite aus Dividenden und Kurssteigerungen des Aktienkurses auch wenn nicht fix verträglich vereinbart
    \end{itemize}
\end{itemize}

\subsubsection{Weighted Average Cost of Capital (WACC)}
\begin{definition}
    [WACC] dient zur Diskontierung erwarteter Gewinne/CashFlows. Besteht aus dem Quotient der erwarteten Renditen von FK-Eigentümer (FK * zinssatz) und EK-Eigentümer (EK * zinssatz) in Bezug auf das Gesamtekapital. Somit muss für jede Einheit gewonnenes Kapital so viel wieder an Eigentümer als Rendite wieder gezahlt werden
    \begin{equation}
        WACC = \frac{EK * k_{EK} + FK * k_{FK} }{GK}
    \end{equation}
\end{definition}

Optimierung der Kapitalstruktur hat das Ziel die Kapitalkosten so niedrig wie möglich zu halten.

\subsection{Ertragswertmethode}

\begin{definition}
    [Ertragswertmethode] Gewinne werden als Diskontierten Cashflows mit Bezug auf die Zukunft berechnet
\end{definition}

\subsubsection{Vorgehen}
\begin{itemize}
    \item Ertragswert (Present Value des Unternehmens) als ewige Rente abstrahiert 
    \item Brutto/Nettoansatz
    \begin{itemize}
        \item Ertragswert netto direkt: $\frac{Reingewinn}{k_{EK}} = Eigenkapital$ 
    \end{itemize}
    \item Ertragswertbrutto: $\frac{Ertragswert + FKZinsen}{WACC} = Gesamtkapital = FK + EK $
    \item Ertragswertnetto: $\frac{Ertragswert + FKZinsen}{WACC} - FK = EK$
\end{itemize}

\begin{definition}
    [Goodwill] Differenz zwischen Ertragswert und Substanzwert des Unternehmens. Sämtliche \textit{immateriellen} Werte die nicht bilanzierbar sind.
\end{definition}

\subsection{Discounted Cashflow Methode (DCF)}



\begin{definition}
    [DCF] Diskontierung aller zukünftigen Free-Cash-FLows nach Abzug von Investitionssumme. FreeCashFlow ist jener Cashflow, der verfügbar ist nach allen notwendigen Erhaltungsnvestitionen. 
\end{definition}

\subsubsection{Berechnung einzelner FCFs}
\begin{itemize}
    \item Gewinn vor Zinsen (Reingewinn + Zinsen)
    \item + Abschreibungen
    \item - Veränderung des operativen Nettoumlaufvermögens
    \item - Investitionen ins Anlagevermögen (wird abgezogen weil nicht \textit{frei} Cashflows sind, sondern notwendig für erhalten der Maschinen/Reinvestition in Kosten)
\end{itemize}

\subsubsection{Vorgang}
\begin{enumerate}
    \item Freie Cashflows berechnen 
    \item FCFs diskontieren 
    \item ewige FCFs berechnen RBF 
    \item Diskontieren der ewigen Rente mit dem Jahr ab dem es beginnt - 1 rechnen
    \item Unternehmenswertbrutto = Summe diskontierten FCFs + diskontierte ewige Rente 
    \item Unternehmenswertnetto (EK) = Brutto - FK Wert
\end{enumerate}